\documentclass[../main.tex]{subfiles}
\begin{document}

\section{Prediction}

$\Ber{T}(x)$ is \emph{correlated} with the latent value $x$, and thus provides information about the latent value. For instance, it allows one to predict $x$ given $\Ber{T}(x)$ better than if no observations were given whatsoever, assuming we know the error rate $\varepsilon(x)$ is better than a random guess, e.g., in the case of $B_{\Bool}$, $\varepsilon(x) < 0.5$.
If we have no prior information about the latent variable, the best (maximum likelihood) estimate of its value is the observation $\Ber{T}(x)$.

However, because $\varepsilon(x)$ is normally known a priori, we can estimate the probability that the latent variable is $X = x$ using Bayes' rule:
\begin{equation}
\Pr\{X = x | \Ber{x} = x\} \propto \Pr\{\Ber{x} = x | X = x\} \Pr\{X = x\},
\end{equation}
where $\Pr\{X = x\}$ is the prior probability that $X = x$ and $\Pr\{\Ber{x} = x | X = x\}$ is the probability that $\Ber{x} = x$ given that $X = x$ (the probability that the observation is correct---the true positive rate $\tau$ or true negative rate $\nu$).

In the first-order model, by Bayes law,
\begin{equation}
\Pr\{x | \Ber{x}\} \propto \tau \Pr\{X = x\},
\end{equation}
where $\tau = 1 - \varepsilon$ is the true positive rate. The probability of observing $x$ given that the latent value is $x$ is $\Pr\{\Ber{x} = x | X = x\} = \tau$. The probability of observing $x$ given that the latent value is not $x$ is $\Pr\{\Ber{x} = x | X \neq x\} = \varepsilon$.

Assuming maximum ignorance (maximum entropy) about $x$ (i.e., $X \sim \rm{Bernoulli}(p=0.5)$), the prediction's posterior above simplifies to $\Pr\{X = x | \Ber{x} = x\} = \tau$ after we normalize the probabilities.

Suppose we have $n$ noisy i.i.d. measurements of the latent variable $x$,
\begin{equation}
\left\{ \Ber{x}, \ldots, \Ber{x} \right\},
\end{equation}
in which case the maximum likelihood estimate of $x$ is the majority vote, i.e., the value that occurs most frequently in the observations.\footnote{Assuming the error rate is not larger than $0.5$.} As the number of independent sources of observations goes to infinity, the majority vote converges in probability to $x$,
\begin{equation}
\lim_{n \to \infty} \Pr\{\rm{majority}(\Ber{x}_1,\cdots,\Ber{x}_n ) = x\} = 1,
\end{equation}
where $\Ber{x}_i$ are independent Bernoulli variables (they do not need to be identically distributed). This is a consequence of the law of large numbers.

This is not a typical use-case for the Bernoulli Boolean model, since it will mostly be an analytical result of probabilistic data structures and algorithms that may be framed in the context of a Bernoulli Model.

\section{Inducing Bernoulli Types}

Here, we discuss how to generalize the results.

\subsection{Unit Functions}

We discussed the idea of the Bernoulli Model for \textbf{value types} like $\Bool$. We can think of these Bernoulli Models as Bernoulli approximations of the set of unit types $() \to X$. A unit type has no input, and it maps to a single constant value in the output. There are $|X|$ functions of this type.

If we replace $() \to X$ with $B_X$ or, equivalently, $B_{() \to X}$, we allow for the possibility of errors. The confusion matrix for $B_{() \to \Bool}^{(2)}$ is provided in Table~\ref{tbl:second_order}.

\begin{table}[htbp]
\centering
\caption{Second-order Bernoulli Boolean Model for $() \to \Bool$}
\label{tbl:second_order}
\begin{tabular}{@{}lcc@{}}
\toprule
$x / B_{\Bool}^{(2)}$ & observe $1$ & observe $0$ \\
\midrule
\textbf{latent} $1$ & $\tau = 1-\eta$ & $\eta$ \\
\textbf{latent} $0$ & $\varepsilon$ & $\nu = 1-\varepsilon$ \\
\bottomrule
\end{tabular}
\end{table}

So, we might observe $B_{() \to \Bool}(x) = 1$ and, according to the confusion matrix, the latent value $x$ is $1$ with probability $\tau$ (true positive rate) and $0$ with probability $\varepsilon$ (false positive rate). Likewise, we might observe $B_{() \to \Bool}(x) = 0$ and the latent value $x$ is $1$ with probability $\eta$ (false negative rate) and $0$ with probability $\nu$ (true negative rate).

We see that the Bernoulli Model for $() \to \Bool$ has a maximum order of 2, since $\eta$ and $\varepsilon$ are independent probabilities that fully describe the model. No additional free parameters are possible, since each row must sum to 1 (the total probability theorem), i.e., given a latent value $x$, the probability of observing $1$ or $0$ must sum to $1$ since those are the only two possible outcomes in the Bernoulli Boolean Model.

We can think of this as the asymmetric binary channel, where the probability of an error is different for $1$ and $0$.

A first-order Bernoulli Boolean model is a model where $\epsilon = \varepsilon = \eta$. See the confusion matrix in Table~\ref{tbl:first_order}.

\begin{table}[htbp]
\centering
\caption{First-Order Bernoulli Boolean Model for $() \to \Bool$}
\label{tbl:first_order}
\begin{tabular}{@{}lcc@{}}
\toprule
$x / B_{\Bool}^{(1)}$ & observe $1$ & observe $0$ \\
\midrule
\textbf{latent} $1$ & $\tau =1-\epsilon$ & $\eta = \epsilon$ \\
\textbf{latent} $0$ & $\varepsilon = \epsilon$ & $\nu = 1-\epsilon$ \\
\bottomrule
\end{tabular}
\end{table}

We can see that there is only one free parameter, $\epsilon$, corresponding to the first-order Bernoulli Boolean Model. We can think of this as a binary symmetric channel, where the probability of an error is the same for $1$ and $0$.

For completeness, we can write down the confusion matrix for the zeroth-order model, which is just the standard Boolean model, in Table~\ref{tbl:zeroth_order}.

\begin{table}[htbp]
\centering
\caption{Zeroth-Order Bernoulli Boolean Model for $() \to \Bool$}
\label{tbl:zeroth_order}
\begin{tabular}{@{}lcc@{}}
\toprule
$x / B_{\Bool}^{(0)}$ & observe $1$ & observe $0$ \\
\midrule
\textbf{latent} $1$ & $1$ & $0$ \\
\textbf{latent} $0$ & $0$ & $1$ \\
\bottomrule
\end{tabular}
\end{table}

We see that there are 0 free parameters, and the model is deterministic.

\subsection{Prediction: Boolean Values}

Earlier, we discussed how to predict latent values from Bernoulli Model observations. Let's apply these insights to the first-order Bernoulli Model for $() \to \Bool$, which we can denote as $B_{\Bool}^{(1)}$.

We can use Bayes' rule to compute the probability that the latent value is $X=1$ given that we observed $B_{\Bool}^{(1)} = 1$:
\begin{equation}
\Pr\{X = 1 | \Berr{1}{1} = 1\} =
    \frac{
        \Pr\{\Berr{1}{1} =1 | X=1\} \Pr\{X = 1\}
    }
    {
        \Pr\{\Berr{1}{1} = 1\}
    }
\end{equation}

We know $\Pr\{\Berr{1}{1}=1 | X = 1\}$ from the confusion matrix; it's just the true positive rate, $\tau = 1 - \epsilon$:
\begin{equation}
\Pr\{X=1 | \Berr{1}{1}=1\} = \frac
    {(1-\epsilon) \Pr\{X = 1\}}
    {\Pr\{\Berr{1}{1}=1\}}
\end{equation}

What is $\Pr\{\Berr{1}{1} = 1\}$? By the total probability theorem:
\begin{align}
\Pr\{\Berr{1}{1} = 1\} &= \Pr\{\Berr{1}{1}=1 | X=1\} \Pr\{X = 1\} + \nonumber\\
    &\quad \Pr\{\Berr{1}{1}=1 | X=0\} \Pr\{X = 0\}.
\end{align}

We can substitute in the confusion matrix values:
\begin{equation}
\Pr\{\Berr{1}{1} = 1\} = (1-\epsilon) \Pr\{X = 1\} + \epsilon \Pr\{X = 0\}.
\end{equation}

If we substitute this back into the expression for $\Pr\{X=1 | \Berr{1}{1}=1\}$, divide the numerator and denominator by $\Pr\{X = 1\}$, and then simplify, we get
\begin{equation}
\Pr\{X=1 | \Berr{1}{1}=1\} = \frac
    {1-\epsilon}
    {1-\epsilon (1 - q/(1-q))}
\end{equation}
where $q = \Pr\{X = 0\}$.

Let's evaluate $q = \Pr\{X = 0\}$ at some interesting points:

\begin{enumerate}
\item At $q = 0$, $\Pr\{X=1 | \Berr{1}{1}=1\} = 1$. This makes sense. We know that the latent value $1$ occurs with probability $1-q=1$, therefore whatever we observe, the latent value is $1$.
\item As $q \to 1$, $\Pr\{X=1 |\Berr{1}{1}=1\} \to 0$. This also makes sense; we know that the latent value $0$ occurs with probability $q$ approaching $1$, so no matter what we observe, the latent value is $0$.
\item At $q = 0.5$, $\Pr\{X=1 |\Berr{1}{1}=1\} = 1-\epsilon$. This is the maximum entropy case, where we have no prior information about the latent value and assume maximum ignorance.
\end{enumerate}

\end{document}
